% RankSTA Paper Part 2: Methodology through Conclusion
% This continues from part1.tex
% Insert this content after the Problem Formulation section

% ==============================================================================
% METHODOLOGY
% ==============================================================================
\section{Methodology}
\label{sec:method}

Figure~\ref{fig:pipeline} illustrates our end-to-end pipeline. Given a Verilog netlist, we: (1) parse the netlist and construct a heterogeneous DAG; (2) extract 10-dimensional node features and 3-dimensional edge features; (3) feed the graph into a 3-layer GAT model; (4) obtain violation probabilities for all nodes; (5) rank nodes and select top-K for inspection.

% FIGURE PLACEHOLDER
\begin{figure}[t]
\centering
\framebox[\columnwidth]{%
\parbox{\columnwidth}{%
\centering
\textcolor{blue}{[FIGURE: End-to-End Pipeline]}\\
Verilog $\rightarrow$ Parser $\rightarrow$ Graph Builder $\rightarrow$ Feature Extractor\\
$\rightarrow$ GAT Model $\rightarrow$ Ranker $\rightarrow$ Top-K Output
}}
\caption{Proposed end-to-end pipeline for timing violation prediction.}
\label{fig:pipeline}
\end{figure}

\subsection{Graph Construction}

We parse the Verilog netlist to extract gate instances (type, input/output pins), net connections (driver pin, load pins), and primary I/O declarations. Algorithm~\ref{alg:dag_build} shows our graph construction procedure.

\begin{algorithm}[t]
\caption{Heterogeneous DAG Construction}
\label{alg:dag_build}
\begin{algorithmic}[1]
\REQUIRE Parsed netlist: gates $\mathcal{G}$, nets $\mathcal{N}$, primary I/O
\ENSURE Directed graph $\mathcal{G} = (\mathcal{V}, \mathcal{E})$, node levels $L$
\STATE Initialize $\mathcal{V} \leftarrow \emptyset$, $\mathcal{E}_{\text{net}} \leftarrow \emptyset$, $\mathcal{E}_{\text{cell}} \leftarrow \emptyset$
\FOR{each primary input $p_i$}
    \STATE $\mathcal{V} \leftarrow \mathcal{V} \cup \{p_i\}$
\ENDFOR
\FOR{each gate $g \in \mathcal{G}$}
    \FOR{each pin $p$ of $g$}
        \STATE $\mathcal{V} \leftarrow \mathcal{V} \cup \{p\}$
    \ENDFOR
    \FOR{each input pin $p_{\text{in}}$ of $g$}
        \STATE $\mathcal{E}_{\text{cell}} \leftarrow \mathcal{E}_{\text{cell}} \cup \{(p_{\text{in}}, p_{\text{out}})\}$ 
    \ENDFOR
\ENDFOR
\FOR{each net $n \in \mathcal{N}$}
    \STATE $p_{\text{driver}} \leftarrow$ driver pin of $n$
    \FOR{each load pin $p_{\text{load}}$ of $n$}
        \STATE $\mathcal{E}_{\text{net}} \leftarrow \mathcal{E}_{\text{net}} \cup \{(p_{\text{driver}}, p_{\text{load}})\}$
    \ENDFOR
\ENDFOR
\STATE $\mathcal{E} \leftarrow \mathcal{E}_{\text{net}} \cup \mathcal{E}_{\text{cell}}$
\STATE $L \leftarrow$ TopologicalSort($\mathcal{G}$)
\RETURN $\mathcal{G}, L$
\end{algorithmic}
\end{algorithm}

For each gate, we create nodes for all pins and add \textit{cell edges} from inputs to outputs. For each net, we add \textit{net edges} from driver to all loads. We then perform topological sorting to compute node levels, which serve as features for downstream GNN processing.

\subsection{Feature Engineering}
\label{sec:features}

Table~\ref{tab:features} lists the 10-dimensional node feature vector. These features capture structural properties (fanout, fan-in, topological level), functional properties (gate type, pin type), timing estimates (delay, slew), and positional encodings.

\begin{table}[t]
\centering
\caption{Node Features (10-Dimensional)}
\label{tab:features}
\begin{tabular}{@{}clp{4.5cm}@{}}
\toprule
\textbf{\#} & \textbf{Feature} & \textbf{Rationale} \\
\midrule
1 & Cell Type & Different gates have different delay characteristics \\
2 & Fanout & High fanout $\rightarrow$ increased capacitive load $\rightarrow$ slower \\
3 & Fan-in & Affects input transition time \\
4 & Topo. Level & Critical paths are typically deep \\
5 & Est. Delay & First-order timing approximation \\
6 & Est. Slew & Slew degradation accumulates \\
7 & Pin Type & Outputs more likely to be endpoints \\
8--10 & Positional & Graph structure embeddings \\
\bottomrule
\end{tabular}
\end{table}

All features are normalized using StandardScaler (zero mean, unit variance) fitted on the training set to prevent design-specific bias. For each edge $(u,v)$, we encode: edge type $t \in \{0, 1\}$ (0=net, 1=cell), estimated delay $d_{uv}$, and load capacitance $c_v$ (fanout-based estimate).

\subsection{GNN Architecture}

We employ a 3-layer Graph Attention Network (GAT) with multi-head attention. The architecture is:
\begin{equation}
\begin{aligned}
\text{Layer 1:} \quad & \mathbb{R}^{10} \xrightarrow{\text{GAT}(h=4)} \mathbb{R}^{128} \xrightarrow{\text{ReLU, Dropout}} \\
\text{Layer 2:} \quad & \mathbb{R}^{512} \xrightarrow{\text{GAT}(h=4)} \mathbb{R}^{128} \xrightarrow{\text{ReLU, Dropout}} \\
\text{Layer 3:} \quad & \mathbb{R}^{512} \xrightarrow{\text{GAT}(h=1)} \mathbb{R}^{2}
\end{aligned}
\label{eq:architecture}
\end{equation}
where $h$ denotes the number of attention heads, dropout rate is 0.2. The first two layers use head concatenation (output dimension $128 \times 4 = 512$), while the final layer averages heads to produce class logits.

For node $i$ with neighbors $\mathcal{N}(i)$, the GAT layer computes:
\begin{equation}
\mathbf{h}_i^{(\ell+1)} = \sigma \left( \sum_{j \in \mathcal{N}(i)} \alpha_{ij} \mathbf{W}^{(\ell)} \mathbf{h}_j^{(\ell)} \right)
\label{eq:gat_update}
\end{equation}
where attention coefficients $\alpha_{ij}$ are learned via:
\begin{equation}
\alpha_{ij} = \frac{\exp(\text{LeakyReLU}(\mathbf{a}^T [\mathbf{W}\mathbf{h}_i \| \mathbf{W}\mathbf{h}_j \| \mathbf{e}_{ij}]))}{\sum_{k \in \mathcal{N}(i)} \exp(\text{LeakyReLU}(\mathbf{a}^T [\mathbf{W}\mathbf{h}_i \| \mathbf{W}\mathbf{h}_k \| \mathbf{e}_{ik}]))}
\label{eq:attention}
\end{equation}

The $\mathbf{e}_i$$j$ term incorporates edge features, allowing the model to learn that certain edge types (e.g., high-delay nets) are more timing-critical.

\textbf{Why GAT over GCN/GraphSAGE?}
(1) \textit{Learned edge importance}: Attention mechanism learns which neighbors matter for timing criticality;
(2) \textit{Multi-head diversity}: Different heads can specialize in different timing modes;
(3) \textit{Edge feature integration}: Unlike GCN, GAT naturally handles edge attributes.

\subsection{Training Strategy}

Due to extreme class imbalance (0.6\% violations on average), standard Cross-Entropy Loss leads to model collapse. We use \textit{weighted Cross-Entropy Loss}:
\begin{equation}
\mathcal{L} = -\frac{1}{|\mathcal{M}|} \sum_{v \in \mathcal{M}} w_{y_v} \log p(y_v | \mathbf{x}_v)
\label{eq:loss}
\end{equation}
where class weights are $w_0 = 1.0$ (safe) and $w_1 = 15.0$ (violation), and $\mathcal{M} = \{v : y_v \neq -1\}$ is the set of labeled nodes (only timing endpoints). This heavily penalizes false negatives, prioritizing recall over precision.

We use Adam optimizer with learning rate $10^{-3}$, weight decay $10^{-5}$ (L2 regularization), and gradient clipping (max norm = 1.0). Training is performed for up to 200 epochs with early stopping (patience = 20 epochs) based on validation ROC-AUC.

\subsection{Ranking-Based Prediction Strategy}

\textbf{Motivation:}
Traditional GNN classifiers use a fixed probability threshold (e.g., $p > 0.5$) to determine violations. However, we discovered empirically that this approach \textbf{fails across different circuits} due to:
\begin{enumerate}
    \item \textbf{Model ``Paranoia'':} Training with high class weights ($w_1 = 15$) causes the model to output inflated probabilities to maximize recall. This is intentional (we want high recall), but it means absolute probabilities are not calibrated across designs.
    
    \item \textbf{Per-Circuit Probability Shifts:} Figure~\ref{fig:prob_dist} shows that:
    \begin{itemize}
        \item Clean circuits: False alarms at $p \approx 0.16$
        \item Violating circuits: True violations at $p \approx 0.08$
    \end{itemize}
    No universal threshold works!
\end{enumerate}

\textbf{Our Solution: Ranking Instead of Thresholding.}
Rather than asking ``Is $p_v > \theta$?'', we ask ``Is $v$ among the top-K riskiest nodes?'' This shifts focus from \textit{absolute probability} to \textit{relative risk rank}. Algorithm~\ref{alg:ranking} shows our rank-based prediction procedure.

\begin{algorithm}[t]
\caption{Rank-Based Timing Violation Prediction}
\label{alg:ranking}
\begin{algorithmic}[1]
\REQUIRE Trained model $f_\theta$, circuit graph $\mathcal{G}$, inspection percentage $K\%$
\ENSURE Top-K risky nodes $\mathcal{V}_{\text{risky}}$
\STATE $\mathbf{p} \leftarrow f_\theta(\mathcal{G})$ \COMMENT{Forward pass}
\STATE $\mathbf{p}_{\text{sorted}} \leftarrow \text{argsort}(\mathbf{p}, \text{descending})$ 
\STATE $K \leftarrow \lceil K\% \times |\mathcal{V}| \rceil$
\STATE $\mathcal{V}_{\text{risky}} \leftarrow \mathbf{p}_{\text{sorted}}[0:K]$
\RETURN $\mathcal{V}_{\text{risky}}$
\end{algorithmic}
\end{algorithm}

\textbf{Advantages of Ranking:}
(1) No per-circuit calibration needed;
(2) Predictable inspection budget (always exactly $K$\% nodes);
(3) Exploits high AUC performance ($>$0.95);
(4) Practical deployment: engineers can plan resource allocation.

\subsection{Adaptive K-Selection Algorithm}

Fixed percentages (e.g., always 5\%) are suboptimal because violation densities vary (0.08\% to 1.11\% in our dataset). We propose an \textit{adaptive K-selection} heuristic:
\begin{equation}
K_{\text{adaptive}} = \max\left( \sum_{v \in \mathcal{V}} \mathbb{1}[p_v > \tau], K_{\text{min}} \right)
\label{eq:adaptive_k}
\end{equation}
where $\tau = 0.01$ is a sensitivity threshold and $K_{\text{min}} = 100$ ensures a minimum inspection budget even for ultra-clean designs.

\textbf{Rationale:} This heuristic works because:
\begin{itemize}
    \item \textbf{Violating circuits:} Many nodes exceed $\tau = 0.01$ $\rightarrow$ Large K $\rightarrow$ Catches all violations
    \item \textbf{Clean circuits:} Few nodes exceed $\tau = 0.01$ $\rightarrow$ Small K $\rightarrow$ Minimizes false inspections
\end{itemize}

Algorithm~\ref{alg:adaptive} shows the complete procedure.

\begin{algorithm}[t]
\caption{Adaptive K-Selection}
\label{alg:adaptive}
\begin{algorithmic}[1]
\REQUIRE Violation probabilities $\mathbf{p}$, sensitivity $\tau$, minimum K $K_{\text{min}}$
\ENSURE Adaptive inspection budget $K$
\STATE $K \leftarrow \sum_{i=1}^{|\mathcal{V}|} \mathbb{1}[p_i > \tau]$
\STATE $K \leftarrow \max(K, K_{\text{min}})$
\RETURN $K$
\end{algorithmic}
\end{algorithm}

% ==============================================================================
% EXPERIMENTAL SETUP
% ==============================================================================
\section{Experimental Setup}
\label{sec:setup}

\subsection{Dataset}

We use the TimingPredict dataset~\cite{guo2022timingpredict}, consisting of 21 open-source VLSI designs (ISCAS, IWLS, OpenCores benchmarks) synthesized with OpenROAD and Sky130 PDK~\cite{sky130pdk}. Each design is synthesized to a gate-level netlist with timing constraints.

Timing labels are extracted via OpenSTA using the following protocol:
(1) Read synthesized netlist and SDC constraints;
(2) Report all timing endpoints: \texttt{all\_registers -data\_pins};
(3) Extract slack: $\text{Slack}_v = \texttt{sta::pin\_slack}(v, \text{max})$;
(4) Binarize: $y_v = 1$ if $\text{Slack}_v < 0$, else $y_v = 0$.

Table~\ref{tab:dataset} shows dataset statistics.

\begin{table}[t]
\centering
\caption{Dataset Statistics}
\label{tab:dataset}
\begin{tabular}{@{}lrrrrr@{}}
\toprule
\textbf{Split} & \textbf{\# Designs} & \textbf{Nodes} & \textbf{Viol.} & \textbf{Rate} & \textbf{Size} \\
\midrule
Train & 12 & 287,543 & 1,784 & 0.62\% & 10K--50K \\
Val & 5 & 91,557 & 621 & 0.68\% & 15K--25K \\
Test & 4 & 221,877 & 1,347 & 0.61\% & 35K--70K \\
\midrule
\textbf{Total} & \textbf{21} & \textbf{600,977} & \textbf{3,752} & \textbf{0.62\%} & \textbf{10K--70K} \\
\bottomrule
\end{tabular}
\end{table}

To ensure balanced violation distribution, we manually assign designs to splits: training includes high-violation designs (e.g., aes256, jpeg\_encoder) plus regularization designs (low-violation), validation has diverse rates (0.1--1\%), and test contains unseen designs with varying densities. This prevents the model from simply memorizing design-specific patterns.

\subsection{Evaluation Metrics}

We use the following metrics:

\textbf{ROC-AUC:} Measures the model's ability to rank violations higher than clean nodes across all thresholds. A score of 1.0 is perfect ranking; 0.5 is random.

\textbf{PR-AUC:} Area Under Precision-Recall Curve, more appropriate for imbalanced datasets than ROC-AUC. Focuses on the positive class.

\textbf{Recall@K\%:} Fraction of true violations caught by inspecting top K\% nodes:
\begin{equation}
\text{Recall@K\%} = \frac{|\{v : v \in \mathcal{V}_{\text{risky}}^{K\%} \wedge y_v = 1\}|}{|\{v : y_v = 1\}|}
\label{eq:recall_k}
\end{equation}

\textbf{Precision@K\%:} Fraction of inspected nodes that are true violations:
\begin{equation}
\text{Precision@K\%} = \frac{|\{v : v \in \mathcal{V}_{\text{risky}}^{K\%} \wedge y_v = 1\}|}{|\mathcal{V}_{\text{risky}}^{K\%}|}
\label{eq:prec_k}
\end{equation}

\subsection{Baselines}

We compare against:
(1) \textbf{Random Baseline}: Randomly rank nodes;
(2) \textbf{Degree Centrality}: Rank by fanout (heuristic: high fanout $\rightarrow$ critical);
(3) \textbf{Longest Path Heuristic}: Rank by topological level (deep nodes $\rightarrow$ critical);
(4) \textbf{GCN}: Replace GAT with Graph Convolutional Network;
(5) \textbf{GraphSAGE}: Replace GAT with GraphSAGE (neighbor sampling);
(6) \textbf{Threshold-Based GAT}: Same GAT but with fixed threshold (0.5, 0.58).

\subsection{Implementation Details}

\textbf{Framework:} PyTorch 2.0~\cite{pytorch_geometric}, PyTorch Geometric 2.3  
\textbf{Hardware:} NVIDIA RTX 3090 (24GB VRAM), Intel Xeon CPU  
\textbf{Training Time:} $\sim$2 hours for 200 epochs (early stopping typically at epoch 80--120)  
\textbf{Inference Time:} $<$50ms per design on GPU, $<$100ms on CPU  
\textbf{Hyperparameter Selection:} Grid search on validation set for learning rate $\{10^{-4}, 10^{-3}, 10^{-2}\}$, hidden dimension $\{64, 128, 256\}$, attention heads $\{2, 4, 8\}$.

\textbf{Code Availability:} All code, pre-trained models, and datasets are available at: \url{https://github.com/Niccc8/gnn-static-timing-analysis-predictor} (MIT License).

% ==============================================================================
% RESULTS AND ANALYSIS
% ==============================================================================
\section{Results and Analysis}
\label{sec:results}

\subsection{Overall Performance}

Table~\ref{tab:results} summarizes the performance of RankSTA and baseline methods on the test set. Our approach achieves \textbf{0.96 ROC-AUC} and \textbf{0.92 PR-AUC}, significantly outperforming all baselines.

\begin{table}[t]
\centering
\caption{Test Set Performance (4 Unseen Designs)}
\label{tab:results}
\begin{tabular}{@{}lcccc@{}}
\toprule
\textbf{Method} & \textbf{ROC-AUC} & \textbf{PR-AUC} & \textbf{Recall@5\%} & \textbf{Time} \\
\midrule
Random & 0.50 & 0.006 & 5.1\% & -- \\
Degree & 0.62 & 0.015 & 18.4\% & $<$1ms \\
Longest Path & 0.68 & 0.023 & 24.7\% & $<$1ms \\
GCN & 0.89 & 0.71 & 62.3\% & 45ms \\
GraphSAGE & 0.91 & 0.78 & 68.1\% & 48ms \\
Threshold-GAT (0.5) & 0.96 & 0.91 & 45.2\% & 42ms \\
Threshold-GAT (0.58) & 0.96 & 0.91 & 51.7\% & 42ms \\
\midrule
\textbf{RankSTA (Ours)} & \textbf{0.96} & \textbf{0.92} & \textbf{78.4\%} & \textbf{41ms} \\
\bottomrule
\end{tabular}
\end{table}

\textbf{Key Observations:}
(1) GNN-based methods significantly outperform heuristics, demonstrating the value of learned representations;
(2) GAT's attention mechanism improves ROC-AUC by 5--7 points over GCN/GraphSAGE;
(3) \textbf{Ranking $>$ Threshold}: Threshold-based GAT achieves the same ROC-AUC (0.96) as ours but only 51.7\% recall at 5\% inspection, while our ranking strategy achieves 78.4\% recall---a \textbf{52\% relative improvement}!

\subsection{Ranking Performance at Different K}

Table~\ref{tab:recall_k} shows recall and precision at different top-K percentages. As K increases, recall improves monotonically.

\begin{table}[t]
\centering
\caption{Recall and Precision at Different Top-K}
\label{tab:recall_k}
\begin{tabular}{@{}lrrrrr@{}}
\toprule
\textbf{Top K\%} & \textbf{Nodes} & \textbf{Recall} & \textbf{Prec.} & \textbf{Gain} & \textbf{Use Case} \\
\midrule
1\% & $\sim$2,200 & 27.4\% & 21.1\% & 100$\times$ & Ultra-fast triage \\
3\% & $\sim$6,600 & 60.8\% & 17.1\% & 33$\times$ & Balanced \\
\textbf{5\%} & \textbf{$\sim$11,000} & \textbf{78.4\%} & \textbf{10.6\%} & \textbf{20$\times$} & \textbf{Recommended} \\
10\% & $\sim$22,000 & 92.1\% & 5.8\% & 10$\times$ & High recall \\
20\% & $\sim$44,000 & 99.2\% & 3.1\% & 5$\times$ & Near-complete \\
\bottomrule
\end{tabular}
\end{table}

\textbf{Recommendation:} Top 5\% provides the best recall-cost trade-off, catching 78\% of violations with 20$\times$ efficiency gain over full analysis.

\subsection{Ablation Studies}

\textbf{Number of Layers:} Table~\ref{tab:ablation_layers} shows that 3 layers is optimal. 2 layers underfit (insufficient receptive field), while 4--5 layers overfit and suffer from vanishing gradients.

\begin{table}[t]
\centering
\caption{Ablation Study: Number of Layers}
\label{tab:ablation_layers}
\begin{tabular}{@{}lcccc@{}}
\toprule
\textbf{\# Layers} & \textbf{ROC-AUC} & \textbf{PR-AUC} & \textbf{Recall@5\%} & \textbf{Params} \\
\midrule
2 & 0.91 & 0.83 & 64.2\% & 410K \\
\textbf{3 (Ours)} & \textbf{0.96} & \textbf{0.92} & \textbf{78.4\%} & \textbf{512K} \\
4 & 0.95 & 0.90 & 75.1\% & 650K \\
5 & 0.93 & 0.87 & 70.3\% & 820K \\
\bottomrule
\end{tabular}
\end{table}

\textbf{Attention Heads:} Table~\ref{tab:ablation_heads} shows that 4 heads strike the best balance.

\begin{table}[t]
\centering
\caption{Ablation Study: Attention Heads}
\label{tab:ablation_heads}
\begin{tabular}{@{}lcccc@{}}
\toprule
\textbf{\# Heads} & \textbf{ROC-AUC} & \textbf{PR-AUC} & \textbf{Recall@5\%} & \textbf{Time} \\
\midrule
1 & 0.93 & 0.88 & 72.1\% & 35ms \\
2 & 0.95 & 0.90 & 75.8\% & 38ms \\
\textbf{4 (Ours)} & \textbf{0.96} & \textbf{0.92} & \textbf{78.4\%} & \textbf{41ms} \\
8 & 0.96 & 0.91 & 77.2\% & 52ms \\
\bottomrule
\end{tabular}
\end{table}

\textbf{Class Weights:} Table~\ref{tab:ablation_weights} shows that high class weights (15$\times$) are necessary to prevent model collapse.

\begin{table}[t]
\centering
\caption{Ablation Study: Class Weights}
\label{tab:ablation_weights}
\begin{tabular}{@{}lcccc@{}}
\toprule
\textbf{Weight $(w_1)$} & \textbf{ROC-AUC} & \textbf{PR-AUC} & \textbf{Recall@5\%} & \textbf{Prec.@5\%} \\
\midrule
1.0 & 0.52 & 0.01 & 3.2\% & 0.4\% \\
3.0 & 0.78 & 0.45 & 42.1\% & 5.1\% \\
10.0 & 0.92 & 0.86 & 69.3\% & 8.9\% \\
\textbf{15.0 (Ours)} & \textbf{0.96} & \textbf{0.92} & \textbf{78.4\%} & \textbf{10.6\%} \\
20.0 & 0.96 & 0.91 & 82.1\% & 9.2\% \\
\bottomrule
\end{tabular}
\end{table}

% ==============================================================================
% DISCUSSION
% ==============================================================================
\section{Industrial Use Case and Discussion}
\label{sec:discussion}

\subsection{Engineering Change Order (ECO) Workflow}

\textbf{Scenario:} Late-stage design with timing violations. Full re-optimization (Place $\rightarrow$ Route $\rightarrow$ STA) takes 20+ hours---unacceptable near tape-out.

\textbf{RankSTA-Enabled ECO:}
\begin{enumerate}
    \item \textbf{Predict} ($<$1 sec): RankSTA identifies top 5\% risky nodes
    \item \textbf{Verify} (10 mins): Run targeted STA: \texttt{report\_timing -to [risky\_list]}
    \item \textbf{Fix} (1 hour): Apply surgical fixes (buffer insertion, gate sizing) only to confirmed violations
    \item \textbf{Iterate} (11 mins total): Repeat if needed
\end{enumerate}

\textbf{Speedup:} 20 hours $\rightarrow$ 1--2 hours per iteration (10--20$\times$ faster).

\subsection{Limitations and Future Work}

\textbf{Clean Circuit Overhead:} The model flags $\sim$8\% of nodes in clean circuits (false positives). This is the "paranoia cost" of ensuring safety. Future work: train a design-level classifier to filter clean circuits \textit{before} node ranking.

\textbf{Multi-Corner Analysis:} Current work focuses on single-corner (setup, worst-case). Extending to multi-corner multi-mode (MCMM) requires modifications to handle corner-specific features.

\textbf{Slack Regression:} Binary classification is simpler than regression but provides less information. Future work could explore ordinal regression (binning slack values).

% ==============================================================================
% CONCLUSION
% ==============================================================================
\section{Conclusion}
\label{sec:conclusion}

We have presented \textit{RankSTA}, a heterogeneous graph neural network framework for pre-routing STA prediction. By modeling circuits as dual-edge DAGs and using a GAT-based classifier with a ranking strategy, our approach achieves $\sim$96\% ROC-AUC on diverse benchmarks (versus $\sim$83\% for XGBoost) and provides 5--20$\times$ speedup over traditional STA. The ranking-based prediction strategy outperforms threshold-based methods by 52\% in recall, eliminating per-circuit calibration and ensuring predictable inspection budgets. This enables timely identification of critical timing endpoints during early design stages, greatly accelerating design iterations.

In future work, we plan to extend the model to multi-corner analysis, explore slack regression for finer-grained predictions, and integrate the predictor into placement tools for closed-loop optimization. Our open-source code and dataset will facilitate continued research in ML-based timing analysis.

% ==============================================================================
% REPRODUCIBILITY
% ==============================================================================
\section*{Reproducibility Checklist}
\begin{itemize}
  \item \textbf{Data and code:} Provided at \url{https://github.com/Niccc8/gnn-static-timing-analysis-predictor} (MIT License)
  \item \textbf{Dependencies:} PyTorch 2.0, PyTorch Geometric 2.3, OpenSTA, Python 3.10+
  \item \textbf{Training:} Adam optimizer (LR=$10^{-3}$), 200 epochs, batch size=1 (full-graph), class weights [1.0, 15.0]
  \item \textbf{Hardware:} NVIDIA RTX 3090 GPU (24GB VRAM)
  \item \textbf{Seeds:} All random seeds fixed for reproducibility (seed=42)
  \item \textbf{Metrics:} Exact definitions of ROC-AUC, PR-AUC, Recall@K\% are given in Section~\ref{sec:setup}
\end{itemize}
